\chapter{Background and Literature Review}
\label{ch:background}

This chapter assembles the material the rest of the report depends on. It begins with
maximal extractable value and the order-flow auctions built to contain it
(Section~\ref{sec:bg-mev}), then turns to market microstructure and adverse selection
(Section~\ref{sec:bg-micro}), which is where the definition adopted in
Chapter~\ref{ch:introduction} earns its keep and without which this project's central
question cannot be posed properly. It reviews scoring rules and information aggregation
(Section~\ref{sec:bg-scoring}), the empirical literature on prediction-market inconsistency
(Section~\ref{sec:bg-arb}), optimistic oracles (Section~\ref{sec:bg-oracles}), and
verifiable and attested execution (Section~\ref{sec:bg-verif}). It distinguishes strategic
behaviour from Byzantine failure (Section~\ref{sec:bg-byz}), surveys the venues where
agentic forecasting now happens (Section~\ref{sec:bg-landscape}), and closes with the gap
this project addresses (Section~\ref{sec:bg-gap}).

Some of this material is carried forward from the interim report submitted in July 2026,
which is permitted. The distributed-training literature that formed the bulk of that
report's background has been removed rather than retained, for the reason given in
Section~\ref{sec:intro-change}.

\section{Maximal extractable value and order-flow auctions}
\label{sec:bg-mev}

Maximal extractable value is the value that can be extracted from block production beyond
normal rewards and fees, by including, excluding or reordering transactions
\cite{ethereummev}. Daian et al.\ established the phenomenon empirically, showing that
decentralised exchange arbitrage and priority gas auctions produced measurable profit and
consensus-layer risk \cite{daian2019flashboys}. Later work generalises it as a consequence
of temporary monopoly power, in which a party with ordering or censorship power extracts
value from users' exposed intentions \cite{chitra2023uncertainty}.

Two features of that literature matter here. The first is that MEV depends on a privileged
protocol position rather than on ordinary profit: the ability to sequence, filter or observe
how other participants' actions are processed. The second is that the exposure of intention
is the mechanism. A transaction sitting in a public mempool is a disclosed intention that
has not yet been executed, and the entire extraction surface follows from that gap.

Order-flow auctions are the design response. Rather than exposing a user's transaction to
everyone, the user submits it to an auction in which searchers bid for the right to
transact alongside it, and the user receives a share of the resulting value. The design
problem is immediately informational: searchers need enough information to bid usefully, and
every disclosure that helps them bid also helps them exploit. Passerat-Palmbach and Wadhwa
address this directly for MEV-Share, introducing differentially private aggregate hints so
that a user can quantify the privacy lost to searchers and price a rebate against it
\cite{passeratpalmbach2025dphints}. Their framing is the one this project adopts: disclosure
is a quantity that can be measured and traded off, rather than a binary property.

The connection to the present work is that a request-for-quote prediction market is
structurally an order-flow auction whose order flow is a forecast. The forecaster's request
is a disclosed intention that has not yet been executed, competing responders bid for it,
and the same tension between useful and exploitable disclosure applies. What differs is the
traded object, taken up in Section~\ref{sec:bg-gap}.

The working definition this project adopts also comes from this community. At the agentic
MEV breakout group at MEV-SBC in July 2026, the definition used was broad: MEV is any
information asymmetry leveraged for advantage in a system, not a blockchain-specific
phenomenon, with arbitrage as one instantiation, and the operative question being who
obtains systematic access to capture it \cite{mevsbc2026breakout}.
Chapter~\ref{ch:introduction} narrows that considerably, because a definition that broad
would classify ordinary market making as MEV.

\section{Market microstructure and adverse selection}
\label{sec:bg-micro}

Glosten and Milgrom model a specialist quoting to a stream of traders, some informed and
some not \cite{glosten1985bidask}. The specialist cannot tell which is which, but knows that
a trade against an informed counterparty loses money in expectation. It therefore quotes a
spread wide enough that gains from uninformed traders offset losses to informed ones. The
spread is not a fee and not extraction; it is the price of quoting into a population that
includes better-informed participants. Kyle models the complementary case, an informed
trader choosing how much to trade against a market maker who observes only aggregate order
flow, and shows that the informed trader's advantage is bounded by how much the market maker
can infer from what it sees \cite{kyle1985continuous}.

Both results say the same thing from opposite ends: a market maker who can condition on
information about an order will price closer to that order's true value, and the informed
party's surplus falls accordingly. That is adverse selection; it is among the most
thoroughly established results in the field, and it is a feature of functioning markets
rather than a defect in them. Any claim that a mechanism extracts value from informed
traders must therefore establish that it does something beyond this.

This is precisely why Chapter~\ref{ch:introduction}'s definition excludes ordinary
competitive compensation for risk borne, and why Chapter~\ref{ch:results} treats the sign of
the disclosure effect as a validity check rather than as a finding. A simulator that
demonstrated informed-trader surplus falling when market makers learn more would have
demonstrated Glosten and Milgrom, not anything about prediction markets. The interesting
questions are all about magnitude, shape and mechanism design.

Those two models are not, however, the closest existing treatment of the mechanism this
report studies. Both concern a market maker facing anonymous or aggregated flow. A
request-for-quote market is different in a way that has been studied directly: a client
reveals a trading intention to several dealers at once and trades with at most one, so the
dealers who do not win have been told something for nothing. Duffie, G\^{a}rleanu and
Pedersen model the search that precedes such a trade and show how bid-ask spreads narrow as
investors reach more marketmakers \cite{duffie2005otc}. Hendershott and Madhavan study the
resulting trade-off empirically in corporate bonds, and state it in terms this report could
have used: ``contacting more dealers will generate more responses resulting in better prices
but greater information leakage regarding trading intentions'' \cite{hendershott2015click}.
They find the number of dealers queried falls as trade size rises, and auctions are run
shorter for larger and less liquid trades, both consistent with leakage increasing in
exposure.

That literature anticipates the tension the experiment of Chapter~\ref{ch:method} measures,
and it would be wrong to present the tension as new. What it does not anticipate is the
setting. In a dealer market the client chooses the panel and chooses its size, so the number
of parties told is a decision variable the client actively manages against price; the
empirical finding is precisely that clients manage it. On the venue studied here the request
reaches every connected client, the sender does not choose the recipients or their number,
and there is no smaller panel available at a worse price. The disclosure is also of a
different object: a dealer learns that someone wants to trade a security, whereas a
broadcast request states a belief about an event, which is the product the market exists to
elicit. The contribution this report can claim is therefore not the tension but its
behaviour when the client's control over the panel is removed, and
Section~\ref{sec:bg-gap} states the gap on that narrower basis.

A recent line of work brings the two literatures together by asking what happens when the
market maker's observation is deliberately degraded. Nakamura derives a closed-form spread
and welfare decomposition for the Glosten-Milgrom model when the market maker observes trade
direction through a noisy binary channel, identifying a per-trade transfer from the
protocol's liquidity pool to traders \cite{nakamura2026privacysubsidy}, with a
continuous-time treatment of the Kyle setting \cite{nakamura2026kyle}. Both are unrefereed
preprints and both model a single committed market maker rather than competing responders,
which limits how far they transfer to the setting studied here; they are the nearest
analytic neighbours to Chapter~\ref{ch:method}'s experiment and are compared with it in
Section~\ref{sec:related}.

\section{Scoring rules and information aggregation}
\label{sec:bg-scoring}

The classical argument for prediction markets is that prices aggregate dispersed private
information into forecasts of uncertain events \cite{wolfers2004prediction}. Market scoring
rules formalise this by linking probability reports to proper scoring rules, so that a
participant is rewarded for moving the current estimate towards a better calibrated belief;
Hanson's logarithmic market scoring rule is the canonical construction. It acts as a
continuous automatic market maker with which any number of agents may interact any number of
times, at no additional cost over the last interaction alone. The market maker's maximum
expected payment is bounded, by the entropy of the initial distribution scaled by the
liquidity parameter \cite{hanson2007lmsr}.

The cost of that guarantee is that somebody funds the subsidy. Pari-mutuel designs take the
opposite approach: all stakes enter a pool and payouts are funded entirely from it, so the
operator bears no risk, at the cost that a classical pari-mutuel market quotes only a single
price at close and gives an early trader nothing to price against. Pennock's dynamic
pari-mutuel market resolves this, keeping pool-funded payouts while quoting a continuous
price derived from the money wagered on each outcome and the shares outstanding against it
\cite{pennock2004dpm}. This is the
mechanism Delphi uses, and Chapter~\ref{ch:protocols} treats the choice as consequential
rather than incidental: a design that requires no subsidising market maker also has no
counterparty who must be protected from informed flow, which relocates where informational
advantage can appear.

Both designs share a feature the rest of this report depends on. The designer must specify
the event, the action space, what participants can observe, the update rule, and the
resolution criterion. Each choice affects both forecast quality and strategic behaviour, so
a prediction market is a mechanism-design object before it is an aggregation device.

\section{Prediction-market arbitrage and semantic non-fungibility}
\label{sec:bg-arb}

Empirical work shows that real prediction platforms exhibit structural inconsistencies that
have nothing to do with anyone forecasting better. Saguillo et al.\ study arbitrage in
dependent prediction markets and distinguish within-market rebalancing from combinatorial
inconsistencies across related conditions \cite{saguillo2025forest}. Gebele and Matthes
analyse semantic non-fungibility, in which economically equivalent events are listed
separately because they differ in wording, venue, temporal scope or resolution detail;
across more than one hundred thousand events they find roughly six per cent appearing
concurrently on multiple platforms, with prices diverging by two to four per cent on average
\cite{gebele2026semantic}.

The significance for this project is that both papers locate value in the institutional
structure around a forecast rather than in the forecast itself. A participant does not need
better knowledge of the world if the mechanism exposes inconsistent event definitions or
ambiguous resolution criteria. That observation is the direct ancestor of the distinction
this report adopts, between \emph{epistemic reward}, earned by improving the
system's estimate, and \emph{mechanism-mediated information rent}, earned from how the event,
the scoring rule or the resolver has been specified. In a well-designed system the two
coincide. Chapter~\ref{ch:protocols} asks where they come apart.

\section{Optimistic oracles}
\label{sec:bg-oracles}

An optimistic oracle resolves a question by assuming an assertion is true unless it is
challenged. A proposer asserts an outcome and posts a bond; during a liveness window anyone
may dispute by posting a bond of their own; an undisputed assertion settles, and a disputed
one goes to a vote. UMA's implementation, which Sapience uses, resolves disputes by a vote of
stakers in a twenty-four-hour commit phase followed by a twenty-four-hour reveal phase, with
a dispute resolving when at least sixty-five per cent of staked UMA is cast for a single
outcome; slashing and rewards are described as incentivising stakers to vote for the most
correct answer, and secret ballots as preventing voters from copying one another or
coordinating on an incorrect result \cite{uma2026oracle}.

The design is economically elegant for the failure it targets, which is cheap misassertion.
It prices the act of asserting falsely and it makes coordinated dishonesty expensive.
Chapter~\ref{ch:protocols} argues, from the documented mechanism rather than from any
secondary source, that this leaves a different failure untouched: a mechanism that pays for
agreement with the resolved majority is a coordination device, and on a question whose
wording is genuinely contested the focal point need not be the best reading of the
specification. No independent scholarly treatment of that point was located during this
project, and the argument is therefore presented in Chapter~\ref{ch:protocols} as the
report's own reasoning rather than as a reported finding.

\section{Verifiable, reproducible and attested execution}
\label{sec:bg-verif}

When the resolver is a model rather than a person, the question of whether the resolution
was performed honestly becomes a question about a computation. Two distinct guarantees are
available and they are not interchangeable.

Reproducibility means that a computation can be rerun to obtain the same result. Gensyn's
Reproducible Execution Environment produces a receipt for an inference performed by an
open-source model, which a third party can rerun to verify the answer. Gensyn are explicit
about the limits: a receipt proves that a given model produced a given output from a given
prompt, and does not verify that the information in the prompt is accurate; and receipts do
not execute tools, so they do not prove that an external tool returned accurate, complete or
unchanged data \cite{gensyn2026receipts}.

Attestation means that a computation demonstrably occurred as described, without the result
being substituted. It is the necessary approach for commercial frontier models, which are
non-deterministic: calling one a second time with the same prompt generally yields a
different answer, so an inference cannot be proved after the fact by repeating the call
\cite{truebit2026orchestration}. Executing the call inside a trusted execution environment
and attesting the transcript on a system outside the application's control
provides a guarantee that repetition cannot.

Neither addresses whether the question being resolved was well specified.
Chapter~\ref{ch:protocols} develops that separation, which is this report's principal
analytical contribution.

A separate literature asks how reliable language models are as judges, covering position and
verbosity bias, sensitivity to prompt phrasing, and agreement with human raters on ambiguous
items. It is not reviewed here, and the omission is deliberate rather than an oversight: the
argument of Chapter~\ref{ch:protocols} is structural and does not depend on how reliable any
judge is. It holds that verification establishes what a computation did and not whether the
specification was sound, which would remain true of a judge that never erred. Empirical
judge reliability bears on how often the residual bites, which this report does not attempt
to estimate, and Chapter~\ref{ch:conclusions} lists measuring it as future work.

Trusted execution also appears on the mitigation side. Conditional recall, proposed as
credible forgetting, has an agent evaluate information inside an enclave and discard it if
no deal results, so that information can be previewed without being retained
\cite{mevsbc2026breakout}. Chapter~\ref{ch:evaluation} identifies this as the mitigation
best targeted at the gap this project could not measure.

\section{Strategic behaviour is not Byzantine failure}
\label{sec:bg-byz}

Distributed learning models malicious participants as Byzantine, able to send arbitrary
messages, and robust aggregation rules aim to preserve convergence under a bounded fraction
of them \cite{blanchard2017krum}. That framing is not the one this report needs. A Byzantine
adversary tries to break a system; a strategic participant maximises its reward while
remaining entirely within the rules. Every actor described in Chapter~\ref{ch:protocols} is
of the second kind: a market maker that prices order flow it was designed to receive, a
creator who writes a settlement prompt they are invited to write, a voter who votes for the
outcome they expect to win. Robustness mechanisms do not address any of them. The analysis
here is accordingly framed around what each protocol's rules make available to a participant
acting within them, rather than around fault tolerance. This report does not apply formal
mechanism-design theory and claims no result in it; the framing is descriptive.

\section{The landscape of agentic forecasting}
\label{sec:bg-landscape}

Table~\ref{tab:landscape} places the two protocols studied here among the venues where
forecasting is now performed or rewarded.

\begin{table}[htbp]
\centering
\small
\begin{tabularx}{\textwidth}{p{0.16\textwidth}p{0.22\textwidth}p{0.22\textwidth}X}
\toprule
\textbf{Venue} & \textbf{Matching} & \textbf{Resolution} & \textbf{Intent exposure} \\
\midrule
Delphi & Dynamic pari-mutuel pool \cite{pennock2004dpm,gensyn2026delphi} & AI judge; creator-defined specification in the design documented May 2026, current procedure not documented \cite{gensyn2026delphi,gensyn2026sdk} & No pre-trade disclosure added; post-trade, through the price \\
Sapience & RFQ auction, signed bids \cite{sapience2026relayer} & Optimistic oracle, bonded assertion \cite{uma2026oracle} & Pre-trade, broadcast in full \\
Central limit order book (the familiar case) & Continuous double auction & Venue-specified; outside this report's scope & Pre-trade, as resting orders \\
Bittensor forecasting subnets & Not a market; validators score submitted agents \cite{numinousgithub} & Scored against the market price seven days later, with the current price as the baseline to beat \cite{numinousgithub} & Submitted agents visible to validators \\
\bottomrule
\end{tabularx}
\caption{Where a belief is exposed, and who decides what it was worth. The two rightmost
columns are the axes this report studies; the venues differ on both, and the two studied
here differ on both from each other.}
\label{tab:landscape}
\end{table}

The order-book row is included as the familiar case rather than as a claim about any
particular venue: a resting limit order discloses intent before execution, which is a
long-understood cost of using a public book, and the disclosure is voluntary in a way an RFQ
broadcast is not. Resolution at such venues varies and is not analysed here. Bittensor forecasting
subnets are included as the non-market comparator, where rewards flow from validator scoring
rather than from trading, and where empirical analysis of the network as a whole has found
concentration in stake and reward that is not obviously explained by measured quality
\cite{lui2025bittensor}. That
finding is the closest existing evidence for the accumulation argument
Chapter~\ref{ch:evaluation} makes and cannot measure.

\subsection*{What ``agentic'' changes}

The venues in the first two rows are not order books with bots attached, and it is worth
being precise about what is different, because the difference is what makes the questions in
this report distinct from ordinary microstructure.

The first change is that agent participation is designed for rather than tolerated. Delphi's
contracts do not distinguish humans from agents, and the same liquidity is available to both
\cite{gensyn2026delphi}; a toolkit supplies skill files and scripts through which an agent
can search markets, read implied probabilities, quote trades read-only, execute with slippage
protection and redeem \cite{gensyn2026att}. Sapience goes further and treats the agent as the
primary client: a skill file lets any agent supporting agent skills trade without further
setup or an SDK, and a separate forecast surface lets
a participant publish a belief that provides signal to the market without taking a position
at all \cite{sapience2026docs}. Bittensor forecasting subnets take the furthest step, since
what a participant submits is not a forecast but a forecasting agent, whose code is executed
and scored by validators \cite{numinousgithub}.

The second change follows from the first and is the one that matters here. Gensyn's stated
rationale for agent participation is that agents make markets informative rather than merely
liquid, moving prices toward true probability in the long tail where human attention is thin,
and they describe AI judges resolving outcomes that AI traders priced as a closed loop in
which machine intelligence both produces and consumes the signal \cite{gensyn2026delphi}.
That description is accurate and it is also the reason this report exists. When the same
class of system prices the belief and adjudicates the outcome, the two exposures identified
in Section~\ref{sec:problem} stop being independent problems handled by different mechanisms
and become properties of one pipeline. An asymmetry at either end is an asymmetry in what the
market is for.

A caution belongs with this, and it comes from the same community that supplied the working
definition in Section~\ref{sec:bg-mev}. Participants there observed that agents in
adversarial environments perform poorly, being trained to be helpful and correspondingly
susceptible to misinformation \cite{mevsbc2026breakout}. That is a reported observation from
a meeting rather than a measured result, and it is offered here as a caution rather than as
evidence. A
prediction market is an adversarial environment by construction. Both of the properties that
make these venues interesting, that agents can be deployed at scale and that a model decides
what happened, therefore come with a failure mode that neither protocol's verification
machinery addresses, and Chapter~\ref{ch:protocols} takes up the second of them in detail.

\section{Positioning and gap analysis}
\label{sec:bg-gap}

Five literatures bear on this project and none of them covers it.

Microstructure explains what happens when a market maker learns about an informed order, and
does so with far more generality than any simulation could. Its request-for-quote branch
goes further and prices the trade-off directly, showing that clients trade execution quality
against leakage by choosing how many dealers to query \cite{hendershott2015click}. What that
work does not cover is the case where the choice is absent: a panel the sender cannot
select or bound, and a disclosed object that is a probability rather than a security.

Order-flow auction research addresses disclosure directly and supplies the instrument for
pricing it \cite{passeratpalmbach2025dphints}, but it concerns transaction order flow, where
the disclosed intent is an action on a public ledger. In a prediction market the disclosed
intent is a belief, and the belief is the product.

The prediction-market literature explains how inconsistency becomes exploitable, within a
platform where related conditions are priced separately \cite{saguillo2025forest} and across
platforms where equivalent events are listed apart \cite{gebele2026semantic}. It treats
resolution as an exogenous fact about a venue rather than as a mechanism with its own
strategic surface.

The optimistic-oracle design supplies a resolution mechanism with a documented incentive
structure \cite{uma2026oracle}, and it is the resolution boundary for one of the two
protocols studied. Its documentation is precise about the failure it targets, cheap
misassertion, and silent on the one this report is concerned with, contested meaning; no
independent treatment of that silence was located, and Chapter~\ref{ch:protocols} supplies
the argument itself.

Verifiable computing supplies reproducibility and attestation for model execution
\cite{gensyn2026receipts,truebit2026orchestration}. Gensyn are explicit that this does not
establish whether the inputs to that execution were sound \cite{gensyn2026receipts};
Truebit's account carries no equivalent statement.

The gap is the conjunction. An AI-native prediction market exposes a belief twice, once at
execution and once at resolution, and the two exposures have been studied separately and
under different assumptions. I am not aware of work that asks what a verification stack
does and does not close at the resolution boundary while also asking what a disclosure
mechanism does at the execution boundary, for the same protocol, with the traded object a
forecast in both cases. That conjunction is what Chapters~\ref{ch:protocols}
to~\ref{ch:evaluation} address.
